Es sei

\mavergleichskettedisp
{\vergleichskette
{ Y
}
{ =} { \bigcup_{i \in I} V_i
}
{ } {
}
{ } {
}
{ } {
}
}
{}{}{}
eine Überdeckung mit offenen Teilmenge

\mavergleichskette
{\vergleichskette
{ V_i
}
{ \subseteq }{ Y
}
{  }{
}
{    }{
}
{   }{
}
}
{}{}{.}
Dies bedeutet, dass es offene Mengen

\mavergleichskette
{\vergleichskette
{ U_i
}
{ \subseteq }{ X
}
{  }{
}
{    }{
}
{   }{
}
}
{}{}{}
gibt mit

\mavergleichskette
{\vergleichskette
{ V_i
}
{ = }{ Y \cap U_i
}
{  }{
}
{    }{
}
{   }{
}
}
{}{}{.}
Daher ist

\mavergleichskettedisp
{\vergleichskette
{ Y
}
{ \subseteq} { \bigcup_{i \in I} U_i
}
{ } {
}
{ } {
}
{ } {
}
}
{}{}{.}
Wegen der Abgeschlossenheit von $Y$ in $X$ ist $X \setminus Y$ offen und daher ist

\mavergleichskettedisp
{\vergleichskette
{ X
}
{ =} { Y \cup (X \setminus Y)
}
{ =} { \bigcup_{i \in I} U_i  \cup  (X \setminus Y)
}
{ } {
}
{ } {
}
}
{}{}{}
eine offene Überdeckung von $X$. Wegen der Kompaktheit von $X$ gibt es eine endliche Teilüberdeckung

\mavergleichskettedisp
{\vergleichskette
{ X
}
{ =} { \bigcup_{i \in J} U_i  \cup  (X \setminus Y)
}
{ } {
}
{ } {
}
{ } {}
}
{}{}{.}
Dies bedeutet wiederum

\mavergleichskettedisp
{\vergleichskette
{ Y
}
{ =} { \bigcup_{i \in J} U_i
}
{ } {
}
{ } {
}
{ } {}
}
{}{}{.}

 [[Kategorie:Latexseite]]
